% ======================= 第 三 章 =======================
\chapter{月兔的来历与编制}
\chaptersubtitle{ГЛАВА ТРЕТЬЯ}{СОЗДАННЫЕ РУКИ}
\epigraphbox{“率土之滨，莫非王臣。”}{《诗经·小雅·北山》}

\kaishou{一个不生育的族群要有人干活，于是造了一个会生育的族群。本章要说明的是这件事的两个后果：数量失控，以及名字。}

\noindent
上一章说他们的判断停住了。判断停住有一个直接后果，很难看：他们不会改变自己，也就不会给自己添人口。

月之民不生育，这是转化的直接结果。没有新生儿，就没有下一代劳力，也没有下一代来接手那些必须有人看着的机器。一座需要常年维护的城，靠一批不会死的人撑着，理论上可以撑很久；可他们自己也承认，撑不住的是「没有人做粗活」这一件事。他们给出的办法是造人。

\section{造物的底本}

月兔不是自然演化出来的物种，这一点在他们的分类里写得很清楚：月兔不入「土著」，也不入「月之民」，单列一类，与器物、牲畜、工具并列。

它们的底本是一种繁殖力极强的地面动物。\footnote{底本的物种在档案中只写作「兔」，未给学名。本书不代为补足，理由是补上去就成了另一个物种。}改造的方向有三个：加上能握东西的手，加上能听懂指令的听觉，加上服从。这三样加起来，就是一个能自我补充的劳动力。

设计上最关键的一条，是让它们能生育。这一点在档案里被反复强调，理由是直白的：造物若不能自我补充，用尽即无，一座城撑不了两万年。

于是有了第二件事。能生育，就没有上限。两万年间，它们的数量从最初的一批涨到了一百多万，远远超过制造它们的人。

\begin{table}[htbp]
\centering
\caption{月兔与月之民的分项对照}
\label{tab:rabbit}
\small
\begin{tabular}{@{}lll@{}}
\toprule
项 & 月之民 & 月兔\\
\midrule
生育 & 无 & 有，且不受控\\
衰老 & 无 & 有，寿数约为人的六成\\
变化 & 无 & 有（因此需要净秽）\\
数量 & 十一万九千（撤离时） & 约一百零四万三千（推算）\\
名字 & 有，来源不明 & 无，全部由主人给\\
\bottomrule
\end{tabular}
\tabsource{月之民一栏据卷 I 简十四、简二十；月兔一栏据 `ФВ-1968-14` 工籍残卷推算，方法见本章第四节。}
\end{table}

表 \ref{tab:rabbit} 的倒数第二行是本章的题目。两万年里，这座城有一百多万人在干活，而这一百多万人在自己的档案里只占一个工种名。

\section{捣者的编制与轮值}

工籍残卷 `ФВ-1968-14` 只剩十一叶，记的是某一年的编制。十一叶里能读出四类：捣、负、修、杂。

捣是制药，负是搬运，修是修理，杂是其余。四类的定额都记着，捣一类最细：三千人，日三合，多者赏，少者罚。

三千这个数在卷 I 里已经出现过，那里把它当成「编制」。本章要补的是另一半：三千是当班的编制，不是全部人口。

依据是残卷第 7 叶的一张轮值表。表分十二格，每格标一个班次，每班三日。三千人一班，十二班轮转，一年轮完。按此推算，能做这道工序的月兔约三万六千；再加上老幼与不在此册者，与第四章所推的百万之数对得上。

这一点解释了卷 I 留下的一个疑惑：为什么「捣者三千」与「月兔百万」看起来打架。两个数都对，一个说的是编制，一个说的是种群。

\section{命名权与「手」}

月之民称月兔为「手」。\footnote{「手」在中文里兼指人与功能，两义并用而不辨，这正是本条要说的。}这个称呼在他们的文书里是标准用法，不带贬义，也不带亲昵，就是一个功能名。相当于我们说的「劳力」「人手」。用词本身没有问题，问题在于配套的那条规矩：月兔没有自己的名字，所有的名字都由主人给。

这条规矩在工籍上留下了痕迹。残卷第 3 叶有一栏「名」，填法有两种：一种是一个单字，另一种是空白。空白者占七成。

单字的名字是怎么来的，残卷里有一条注：按入职顺序，取天干、地支、方位、数目为名。也就是说，它们的名字是编号的另一种写法，不是称呼。

笔者见过一份例外。第 9 叶记一名月兔因病停工四十日，页边有人补了一行小字，写的是「阿六，西廊第三班，手稳」。这一行超出了编号的用途，是有人真的记住了它。这是全书材料里唯一一次。

\begin{archive}[机密／СЕКРЕТНО]{档案 № ФВ-1968-14}
工籍残卷，存十一叶，纸质，单面墨书，边缘焦脆。原册页码自第 3 叶起，前二叶已失。

第 3 叶为分类总目，列四类：捣、负、修、杂，各类下注定额与轮值周期。

第 7 叶为轮值表，十二格，每格三日，格内填人数，合计三万六千。

第 9 叶记一名月兔因病停工四十日，页边补一行小字：「阿六，西廊第三班，手稳。」

第 11 叶末行有墨涂，涂痕下可辨「不可使聚」四字。
\end{archive}

\section{材料中的空白}

这一节要写的东西，材料里没有。月兔有没有自己的语言、有没有互相称呼的办法、有没有对这座城的看法，档案一个字也没有。它们的记录全部由主人书写，用途是核算劳力，不是理解对方。

笔者能推的只有一件事，而且是从反面推的：既然那一行「阿六，手稳」会被写下来，说明写字的人与这个对象之间存在过某种具体的关系；而这样的记录全书只有一条，说明这种关系不被制度承认。

至于它们自己怎么想，本书不给答案。这个空白出于材料的选择：写档案的人没有理由去问。

\section{本章结论}

月之民造了一个能生育的族群来替自己劳动，两万年后，这个族群有一百多万人，而制造者只剩十一万九千。

本章要说的是这条曲线：

造物能自我补充，因此数量不受控；数量不受控，因此必须有制度来管；制度给了它们工种、班次、定额，却没有给名字。名字由主人给，编号就是名字。这个设计的后果，在第七章兑现。

本章的自我拆台在这里：一百零四万三千这个数，是从一张轮值表按班次推算的，不是清点记录。全书找不到任何一次对月兔的普查。若不承认这个推算，那么它们的数量在本书里就是未知的，本章的「数量失控」也就不成立。

笔者保留这个推算，因为它与另外两个数对得上：捣者的年产量，与蓬莱药在山中的存量。三个数互相独立，得出的量级一致。非普查材料之间能对上量级，通常不会全错。
